% !TEX TS-program = xelatex
% !TEX encoding = UTF-8 Unicode
% ================================================
%  上海交通大学 研究生工作周报
%  编译: xelatex → bibtex → xelatex → xelatex
% ================================================

\documentclass{sjtuweekly}

% ========== 参考文献库 ==========
\bibfile{bib/database}

% ========== 基本信息 ==========
\member{张\quad 三}
\weeknum{1}


\begin{document}

\makeweeklytitle   % 封面（第1页）

% ================================================
%  正文从第2页开始
% ================================================

\section{一、本周工作概要}

\subsection*{1.1\quad 文献阅读}

本周重点研读了 Tussyadiah 等人关于共享经济中用户评价文本分析的论文
\cite{tussyadiah2015hotels}。该工作使用自然语言处理技术对波特兰地区酒店与
P2P 住宿平台的消费者评论进行对比分析，为本课题提供了重要的方法论参考。

\subsection*{1.2\quad 实验进展}

\begin{itemize}
  \item 完成数据集的清洗与预处理，共计处理样本 12\,000 条
  \item 实现基线模型并在验证集上评估，初步结果如表~\ref{tab:result} 所示
  \item 针对类别不平衡问题尝试了重采样策略
\end{itemize}

\begin{table}[htbp]
  \centering
  \begin{tabular}{lcc}
    \toprule
    方法        & 准确率  & F1 分数 \\
    \midrule
    Baseline    & 0.851  & 0.823 \\
    + Resample  & 0.864  & 0.841 \\
    Ours (v1)   & 0.892  & 0.871 \\
    \bottomrule
  \end{tabular}
  \caption{各方法在验证集上的性能对比}
  \label{tab:result}
\end{table}

\subsection*{1.3\quad 理论分析}

KL 散度是度量两个概率分布差异的基本工具：

\begin{equation}
  D_{\mathrm{KL}}(p \parallel q)
    = \int p(\boldsymbol{x}) \,
      \ln \frac{p(\boldsymbol{x})}{q(\boldsymbol{x})} \,
      \mathrm{d}\boldsymbol{x},
  \label{eq:KL}
\end{equation}

该量在变分推断与生成模型中具有核心地位（式~\ref{eq:KL}）。

% ================================================
\section{二、遇到的问题与讨论}

\begin{enumerate}
  \item \textbf{数据不平衡}：正负样本比例约为 1:8，Focal Loss 在初步实验中
    未能显著提升少数类的召回率，考虑结合过采样策略进一步验证。
  \item \textbf{收敛速度}：模型在前 20 个 epoch 内 loss 下降缓慢，怀疑与学习
    率预热策略的设置有关，下周将进行系统的超参数搜索。
\end{enumerate}

% ================================================
\section{三、下周工作计划}

\begin{itemize}
  \item 完成消融实验的全部配置并提交计算任务
  \item 阅读 Partl 等人的技术写作文献 \cite{partl2016}，为方法部分撰写做准备
  \item 整理实验日志，补充缺失的随机种子记录
\end{itemize}

% ================================================
\printbib

\end{document}
